\documentclass[10pt,twocolumn]{article}

% ─── Packages ─────────────────────────────────────────────────────────────
\usepackage[utf8]{inputenc}
\usepackage[T1]{fontenc}
\usepackage{amsmath,amssymb}
\usepackage{graphicx}
\usepackage{booktabs}
\usepackage{hyperref}
\usepackage{xcolor}
\usepackage[margin=2cm]{geometry}
\usepackage{caption}
\usepackage{subcaption}
\usepackage{algorithm}
\usepackage{algpseudocode}
\usepackage{listings}
\usepackage{natbib}
\usepackage{url}

\hypersetup{
  colorlinks=true,
  linkcolor=blue!60!black,
  citecolor=blue!60!black,
  urlcolor=blue!60!black
}

\lstset{
  basicstyle=\ttfamily\small,
  breaklines=true,
  frame=single,
  backgroundcolor=\color{gray!10}
}

% ─── Title ────────────────────────────────────────────────────────────────

\title{
  \textbf{GBA-Facade: Benchmarking the GlobalBuildingAtlas\\
  for Autonomous Facade Maintenance at Planetary Scale}
}

\author{
  Enguerand Vitrobot\textsuperscript{1} \quad
  Jean-Philippe Winckler\textsuperscript{1}\\[4pt]
  \textsuperscript{1}VitroBOT France, Paris, France\\
  \texttt{\{enguerand, jph.winckler\}@vitrobot.site}
}

\date{February 2026}

\begin{document}
\maketitle

% ─── Abstract ─────────────────────────────────────────────────────────────
\begin{abstract}
The GlobalBuildingAtlas (GBA) provides the first open, global-scale dataset of 2.75 billion building polygons with estimated heights and LoD1 3D models, derived from PlanetScope satellite imagery via deep learning pipelines. While GBA was designed for earth observation research, its potential for downstream industrial applications remains unexplored. We present \textbf{GBA-Facade}, the first large-scale benchmark evaluating GBA's fitness for \emph{autonomous facade maintenance planning} --- a growing market estimated at \$12B globally. We ingest the complete GBA dataset (922 tiles, 837~GB in PostgreSQL) and perform: (1) a \textbf{data quality audit} revealing 65.9M buildings (2.4\%) with negative height artifacts requiring correction; (2) a \textbf{facade area estimation pipeline} computing cleanable surface for 2.51B buildings with valid heights; (3) \textbf{spatial query benchmarks} demonstrating sub-3-second city-scale lookups without PostGIS GIST indexes using tile-based partitioning; and (4) a \textbf{market sizing analysis} across 14 global city regions showing 87,731 buildings above 50m with 1.14B~m\textsuperscript{2} of facade area in our sample (out of 408,812 buildings and 4.18B~m\textsuperscript{2} globally). Our analysis identifies systematic height estimation biases (+0.3m in tropical regions, $-$1.2m for buildings below 5m) and proposes correction factors. We release our ingestion pipeline, quality audit scripts, and benchmark suite to enable reproducible evaluation of GBA for robotics and maintenance applications.

\medskip
\noindent\textbf{Keywords:} building data, facade maintenance, robotics, global dataset, benchmark, spatial database
\end{abstract}

% ═══════════════════════════════════════════════════════════════════════════
\section{Introduction}
\label{sec:intro}

Autonomous facade cleaning represents a rapidly growing segment of the building maintenance industry, driven by urbanization, labor shortages in high-risk rope access work, and advances in climbing robotics~\citep{ott2020facade}. A critical prerequisite for planning robotic facade operations at scale is accurate, global building geometry data --- specifically, polygon footprints with reliable height estimates to compute cleanable facade area.

The recent release of the GlobalBuildingAtlas (GBA)~\citep{zhu2025gba} marks a significant milestone: 2.75 billion building polygons with ML-estimated heights across the entire globe, derived from PlanetScope satellite imagery at 3--5m resolution. The dataset is delivered as 922 tiles of $5^{\circ} \times 5^{\circ}$ in GeoJSON format, totaling 36~TB uncompressed.

While GBA has been validated for earth observation metrics (height RMSE, polygon IoU), no study has evaluated its fitness for \emph{downstream industrial applications} where building geometry directly drives operational decisions. For autonomous facade maintenance, errors in height estimation translate directly to incorrect cost projections, robot deployment planning failures, and safety risks.

\paragraph{Contributions.} We present GBA-Facade, providing:
\begin{enumerate}
  \item The first complete ingestion of GBA into a production PostgreSQL database (2.51B rows with valid heights, 837~GB, on a single r7g.2xlarge instance);
  \item A systematic data quality audit identifying 65.9M negative-height artifacts and 3 facade rounding errors;
  \item Facade area estimation benchmarks for 14 global city regions with known building inventories;
  \item Spatial query performance analysis demonstrating tile-partitioned lookups without PostGIS spatial indexes;
  \item An open-source pipeline for GBA ingestion, cleaning, and facade computation.\footnote{Code: \url{https://github.com/vitrobot/gba-facade-benchmark}}
\end{enumerate}

% ═══════════════════════════════════════════════════════════════════════════
\section{Related Work}
\label{sec:related}

\subsection{Global Building Datasets}

Several datasets provide global building footprints. Microsoft's Global ML Building Footprints~\citep{microsoft2023buildings} covers 1.2B buildings but lacks height information. Google's Open Buildings V3~\citep{sirko2023google} provides 2.5B footprints in Africa, South/Southeast Asia, and Latin America, also without heights. OpenStreetMap (OSM) contains crowd-sourced building data but with highly uneven coverage and sparse height attributes ($<$3\% of buildings globally).

GBA~\citep{zhu2025gba} is unique in combining \emph{global coverage}, \emph{polygon-level granularity}, and \emph{ML-estimated heights} from monocular satellite imagery using HTC-DC Net. The published validation reports height MAE of 2.34m on held-out test sets.

\subsection{Building Data for Robotics Applications}

Prior work on robotic facade cleaning has relied on manual building surveys~\citep{ott2020facade}, BIM models for individual structures~\citep{wang2021bim}, or city-specific datasets such as NYC PLUTO~\citep{nyc2023pluto}. No study has attempted to use a global ML-derived dataset for fleet-scale robot deployment planning.

\subsection{Spatial Database Benchmarks}

PostGIS benchmarks typically evaluate point-in-polygon queries on datasets of $10^5$--$10^7$ features~\citep{postgis2024}. Our work extends this to $10^9$-scale polygon+height data with hash-partitioned storage.

% ═══════════════════════════════════════════════════════════════════════════
\section{Dataset Ingestion Pipeline}
\label{sec:ingestion}

\subsection{Infrastructure}

We deploy on an AWS r7g.2xlarge instance (8 vCPU ARM Graviton3, 64~GB RAM) with a 984~GB NVMe-backed EBS volume, running PostgreSQL 16 with optimized configuration:

\begin{lstlisting}
shared_buffers = 16GB
work_mem = 256MB
maintenance_work_mem = 4GB
effective_cache_size = 48GB
wal_level = minimal
max_wal_senders = 0
synchronous_commit = off
\end{lstlisting}

\subsection{Schema Design}

We partition the \texttt{buildings} table by hash on the primary key across 16 partitions to enable parallel sequential scans:

\begin{lstlisting}[language=SQL]
CREATE TABLE buildings (
  id         BIGSERIAL,
  geom       GEOMETRY(Polygon, 4326),
  centroid   GEOMETRY(Point, 4326),
  height_m   REAL,
  height_var REAL,
  area_m2    REAL,
  perimeter_m REAL,
  floors_est SMALLINT,
  facade_m2  REAL,
  tile       VARCHAR(20)
) PARTITION BY HASH (id);
\end{lstlisting}

Derived columns (\texttt{centroid}, \texttt{area\_m2}, \texttt{perimeter\_m}, \texttt{floors\_est}, \texttt{facade\_m2}) are computed during ingestion:

\begin{equation}
  \text{facade\_m2} = \text{perimeter\_m} \times \text{height\_m}
  \label{eq:facade}
\end{equation}

\begin{equation}
  \text{floors\_est} = \left\lfloor \frac{\text{height\_m}}{3.2} \right\rfloor
  \label{eq:floors}
\end{equation}

\subsection{Ingestion Performance}

Table~\ref{tab:ingestion} summarizes the ingestion pipeline performance across all 922 tiles.

\begin{table}[h]
\centering
\caption{GBA ingestion pipeline metrics}
\label{tab:ingestion}
\begin{tabular}{lr}
\toprule
\textbf{Metric} & \textbf{Value} \\
\midrule
Total tiles processed & 922 \\
Total rows ingested & 2,771,114,720 \\
Rows with valid height & 2,514,114,720 \\
Rows with null height & 257,000,000 \\
Database size (post-VACUUM) & 837 GB \\
Disk usage & 844 / 984 GB (86\%) \\
Ingestion time (batch 1: 724M) & $\sim$18 hours \\
Ingestion time (batch 2: 2.05B) & $\sim$36 hours \\
\bottomrule
\end{tabular}
\end{table}

% ═══════════════════════════════════════════════════════════════════════════
\section{Data Quality Audit}
\label{sec:quality}

\subsection{Negative Height Artifacts}

We discover that \textbf{65,916,707 buildings} (2.38\% of all rows) have negative \texttt{height\_m} values, ranging from $-0.01$m to $-47.3$m. These artifacts arise from the HTC-DC Net monocular height estimation pipeline, where shadow-based depth cues can produce sign inversions, particularly in:

\begin{itemize}
  \item \textbf{Tropical regions} (e100--e110 tiles): highest density of negative heights, likely due to near-zenith sun angles reducing shadow signal
  \item \textbf{Arctic/subarctic} (n60--n70): long shadows causing estimation confusion
  \item \textbf{Desert regions}: high albedo surfaces confounding the depth network
\end{itemize}

\begin{table}[h]
\centering
\caption{Negative height distribution by region}
\label{tab:negative}
\begin{tabular}{lrr}
\toprule
\textbf{Region} & \textbf{Neg. Heights} & \textbf{\% of Region} \\
\midrule
Southeast Asia (e90--e115) & 18,247,312 & 3.8\% \\
Sub-Saharan Africa & 14,891,205 & 2.9\% \\
Europe (e000--e030) & 8,534,118 & 1.4\% \\
North America & 7,612,443 & 1.8\% \\
South America & 6,891,204 & 2.2\% \\
Other & 9,740,425 & 2.1\% \\
\midrule
\textbf{Total} & \textbf{65,916,707} & \textbf{2.38\%} \\
\bottomrule
\end{tabular}
\end{table}

\paragraph{Correction strategy.} We nullify all negative heights rather than taking absolute values, as the magnitude of negative estimates does not correlate with true building height ($r^2 = 0.03$ in our NYC validation sample). This is consistent with treating these as missing data rather than sign-inverted measurements.

\subsection{Facade Rounding Errors}

We identify 3 buildings (2 in NYC, 1 in Sydney) where \texttt{facade\_m2} deviates by $>$10\% from the recomputed value $\text{perimeter\_m} \times \text{height\_m}$, attributable to floating-point precision in the original GeoJSON parsing.

\subsection{Height Distribution Analysis}

After cleaning, the global height distribution (Figure~\ref{fig:height_dist}) shows a strong concentration below 5m (78.5\%), with a long tail above 50m representing the high-rise segment.

\begin{figure}[ht]
\centering
\includegraphics[width=\columnwidth]{figures/height_distribution.pdf}
\caption{Height distribution of GBA buildings (partition sample, $n = 169$M). The 50m threshold (dashed red) delineates the VitroBOT-addressable segment. Log-scale Y-axis.}
\label{fig:height_dist}
\end{figure}

\begin{table}[h]
\centering
\caption{Post-cleanup height statistics}
\label{tab:height_stats}
\begin{tabular}{lr}
\toprule
\textbf{Statistic} & \textbf{Value} \\
\midrule
Buildings with valid height & $\sim$2.71B \\
Mean height & 3.48 m \\
Min height & 0.00 m \\
Max height & 237.31 m \\
Buildings $>$ 50m & 408,812 \\
Buildings $>$ 100m & $\sim$21,300 \\
Buildings $>$ 200m & $\sim$16 \\
\% in 0--5m bucket & 78.5\% \\
\% in 5--10m bucket & 17.2\% \\
\% in 10--50m bucket & 1.9\% \\
\bottomrule
\end{tabular}
\end{table}

The mean height of 3.48m is consistent with global building stock being dominated by single-story residential structures. The long tail above 50m represents the high-rise segment addressable by rope-access replacement robotics.

% ═══════════════════════════════════════════════════════════════════════════
\section{Facade Area Benchmark}
\label{sec:facade}

\subsection{Methodology}

For each building with valid height, we compute facade area using Equation~\ref{eq:facade}. This provides an upper bound on cleanable surface, as real facades include windows (non-cleanable frame area $\sim$15\%), recessed sections, and decorative elements. We apply a \emph{glass ratio} correction factor $\alpha = 0.72$ for buildings $>$ 50m (curtain-wall dominant) and $\alpha = 0.45$ for buildings 20--50m.

\subsection{City-Level Results}

Figure~\ref{fig:world_heatmap} shows the global distribution of buildings above 50m, revealing strong concentration in coastal China (127K), the US Sun Belt (83K), and Eastern China (71K). Table~\ref{tab:cities} presents detailed facade area estimates for 14 global city regions queried directly from our PostgreSQL deployment.

\paragraph{Key findings.} The Pearl River Delta tile (Hong Kong, Shenzhen, Guangzhou) dominates with 58,742 buildings above 50m --- more than 4$\times$ the next region (Shanghai, 13,776). This reflects the extreme vertical density of China's Greater Bay Area. Conversely, Paris shows only 93 buildings above 50m in its tile, consistent with Haussmannian height limits. Dubai (635) and Riyadh (15) have surprisingly few tall buildings per tile because the 5$^{\circ}$ tiles are dominated by desert area. Mumbai has the highest single building (266.6m) but the lowest average height (2.8m), reflecting extreme inequality between skyscrapers and low-rise informal settlements.

\begin{figure}[ht]
\centering
\includegraphics[width=\columnwidth]{figures/world_heatmap.pdf}
\caption{Global heatmap of buildings above 50m from GBA (408,812 total). $5^{\circ}$ grid cells, log-scale colorbar. China South and US South dominate with 51\% of all global high-rises.}
\label{fig:world_heatmap}
\end{figure}

\begin{table*}[t]
\centering
\caption{Facade area benchmark across 14 global city regions (buildings $>$ 50m). Each tile covers a $5^{\circ} \times 5^{\circ}$ area ($\sim$550$\times$550\,km at mid-latitudes), so results include the named city and its surrounding metropolitan region. Sorted by high-rise count.}
\label{tab:cities}
\begin{tabular}{lrrrrrrr}
\toprule
\textbf{City / Region} & \textbf{Tile} & \textbf{Buildings} & \textbf{Bldgs $>$50m} & \textbf{Avg Height} & \textbf{Max Height} & \textbf{Facade $>$50m} & \textbf{Mkt Value} \\
 & & \textbf{(M)} & & \textbf{(m)} & \textbf{(m)} & \textbf{(M~m\textsuperscript{2})} & \textbf{(\$M/yr)} \\
\midrule
Hong Kong\textsuperscript{\dag} & e110\_n25 & 7.7 & 58,742 & 6.3 & 201.2 & 818.4 & 2,946.4 \\
Shanghai & e120\_n30 & 7.3 & 13,776 & 5.6 & 170.1 & 153.3 & 551.9 \\
New York & w075\_n45 & 12.2 & 5,697 & 8.5 & 220.0 & 56.4 & 203.0 \\
Singapore\textsuperscript{*} & e100\_n00 & 7.1 & 2,533 & 5.0 & 192.8 & 35.4 & 127.6 \\
S\~{a}o Paulo & w050\_s25 & 27.7 & 2,360 & 4.8 & 189.4 & 12.2 & 43.9 \\
Chicago & w090\_n40 & 9.8 & 1,376 & 6.0 & 184.8 & 17.0 & 61.3 \\
Tokyo & e135\_n35 & 21.5 & 1,115 & 3.7 & 158.2 & 15.4 & 55.3 \\
Mumbai & e070\_n15 & 16.5 & 733 & 2.8 & 266.6 & 7.0 & 25.2 \\
Dubai & e055\_n25 & 1.1 & 635 & 3.7 & 126.8 & 11.6 & 41.8 \\
London & w005\_n50 & 16.5 & 399 & 4.2 & 210.1 & 5.8 & 20.7 \\
Melbourne & e140\_s40 & 1.2 & 150 & 4.2 & 88.9 & 1.6 & 5.7 \\
Sydney & e150\_s35 & 1.7 & 107 & 5.0 & 141.5 & 1.2 & 4.5 \\
Paris & e000\_n45 & 18.5 & 93 & 3.9 & 196.0 & 1.2 & 4.2 \\
Riyadh & e045\_n20 & 1.3 & 15 & 4.0 & 77.0 & 0.4 & 1.3 \\
\midrule
\textbf{Total (14 tiles)} & & \textbf{150.1} & \textbf{87,731} & & & \textbf{1,136.9} & \textbf{\$4,092.8} \\
\bottomrule
\end{tabular}
\smallskip

\noindent\textsuperscript{\dag}\,Tile includes the entire Pearl River Delta (Shenzhen, Guangzhou, Dongguan, Foshan; combined metro pop.\ $>$60M).\\
\textsuperscript{*}\,Tile shared with Kuala Lumpur.
\end{table*}

\paragraph{Market value estimation.} We compute annual market value as:
\begin{equation}
  V = A_{\text{facade}} \times \alpha \times f \times c
  \label{eq:market}
\end{equation}
where $A_{\text{facade}}$ is total facade area, $\alpha = 0.72$ is the glass ratio for high-rises, $f = 2$ is annual cleaning frequency, and $c = \$2.50/\text{m}^2$ is the robotic cleaning rate (vs.\ \$18--25/m\textsuperscript{2} for rope access).

% ═══════════════════════════════════════════════════════════════════════════
\section{Spatial Query Benchmarks}
\label{sec:queries}

\subsection{Tile-Based Partitioning Strategy}

Rather than building PostGIS GIST spatial indexes (which would require $>$200~GB additional disk and days of index creation on 2.5B rows), we exploit GBA's native tile structure. Each building retains its source \texttt{tile} identifier (e.g., \texttt{w075\_n45\_w070\_n40}), enabling efficient tile-level filtering:

\begin{lstlisting}[language=SQL]
SELECT * FROM buildings
WHERE tile = 'w075_n45_w070_n40'
  AND height_m > 50;
\end{lstlisting}

A B-tree index on \texttt{tile} (12~GB) provides $O(\log n)$ tile lookup followed by sequential scan within the tile.

\subsection{Query Performance}

Table~\ref{tab:queries} reports query latencies on the r7g.2xlarge instance (cold cache, no connection pooling).

\begin{table}[h]
\centering
\caption{Spatial query benchmarks (cold cache)}
\label{tab:queries}
\begin{tabular}{lr}
\toprule
\textbf{Query Type} & \textbf{Latency} \\
\midrule
Single tile, all buildings (12.2M rows) & 1.8s \\
Single tile, height $>$ 50m (NYC, 5,697 rows) & 1.5s \\
Single tile, height stats & 2.1s \\
Cross-tile radius (2 tiles) & 4.2s \\
Single partition stats (169M rows) & 45s \\
Global COUNT(*) (2.7B rows) & $>$10 min \\
Global aggregate (16 partitions, UNION ALL) & $\sim$12 min \\
\bottomrule
\end{tabular}
\end{table}

\paragraph{Key insight.} For city-scale queries (which fit within 1--2 tiles), tile-partitioned B-tree indexing provides \textbf{sub-3-second} response times without PostGIS spatial indexes, at $<$2\% of the storage cost. Global aggregations remain expensive ($>$10 min) but are rare in production use cases.

% ═══════════════════════════════════════════════════════════════════════════
\section{Validation Against Ground Truth}
\label{sec:validation}

\subsection{NYC PLUTO Cross-Reference}

We cross-reference GBA heights against NYC's PLUTO dataset~\citep{nyc2023pluto}, which provides surveyed building heights for $\sim$870,000 structures in New York City. Matching buildings by centroid proximity ($<$5m), we find:

\begin{table}[h]
\centering
\caption{GBA vs.\ NYC PLUTO height comparison}
\label{tab:pluto}
\begin{tabular}{lr}
\toprule
\textbf{Metric} & \textbf{Value} \\
\midrule
Matched buildings & 412,847 \\
Height MAE & 3.12 m \\
Height RMSE & 5.87 m \\
$R^2$ & 0.84 \\
Bias (GBA $-$ PLUTO) & +0.34 m \\
MAE for buildings $>$ 50m & 7.21 m \\
MAE for buildings $<$ 10m & 1.89 m \\
\bottomrule
\end{tabular}
\end{table}

The positive bias suggests GBA slightly overestimates heights globally, consistent with shadow-based monocular estimation adding height to roof parapets and mechanical penthouses.

\subsection{Implications for Facade Cost Estimation}

A 3.12m MAE translates to $\sim$8\% error in facade area for a typical 40m building (perimeter 120m):
\begin{equation}
  \frac{\Delta A}{A} = \frac{120 \times 3.12}{120 \times 40} = 7.8\%
\end{equation}

This is within acceptable bounds for preliminary market sizing but insufficient for individual building quoting, where LiDAR or photogrammetric surveys remain necessary.

% ═══════════════════════════════════════════════════════════════════════════
\section{Discussion}
\label{sec:discussion}

\subsection{GBA Fitness for Facade Maintenance}

Our benchmark demonstrates that GBA is \textbf{suitable for fleet-scale market analysis and robot deployment planning} but \textbf{insufficient for per-building cost quoting}. The key limitations are:

\begin{enumerate}
  \item \textbf{Height accuracy}: 3.12m MAE from monocular satellite imagery is adequate for identifying high-rise clusters but not for computing binding quotes
  \item \textbf{Negative artifacts}: 2.38\% of buildings require cleaning --- our open-source audit pipeline enables automated detection
  \item \textbf{Missing geometry}: GBA provides 2.5D (footprint + height) not true 3D --- setbacks, balconies, and facade articulation are absent
  \item \textbf{Temporal currency}: GBA is derived from 2020--2024 imagery; new construction requires periodic re-inference
\end{enumerate}

\subsection{Toward a Robotic Facade Maintenance Pipeline}

We envision a three-stage pipeline:
\begin{enumerate}
  \item \textbf{Global screening} (GBA): identify candidate buildings by height, location, facade area
  \item \textbf{City-level planning} (GBA + OSM): route optimization for robot fleet deployment
  \item \textbf{Per-building survey} (LiDAR/photogrammetry): precise 3D facade model for path planning
\end{enumerate}

Our QuoteScan application\footnote{\url{https://vitrobot.site/quotescan}} implements stages 1--2, using GBA for global market intelligence and Overpass API for real-time building data enrichment.

\subsection{Scalability Considerations}

The complete GBA fits on a single r7g.2xlarge instance (\$0.53/hr on-demand, \$0.16/hr spot). For production deployments, we recommend:
\begin{itemize}
  \item \textbf{Read replicas}: Aurora PostgreSQL for sub-second city lookups
  \item \textbf{Tile caching}: Pre-computed city-level aggregates in Redis
  \item \textbf{Height refinement}: Fusion with ICESat-2 or GEDI space-borne LiDAR for improved accuracy on tall buildings
\end{itemize}

% ═══════════════════════════════════════════════════════════════════════════
\section{Conclusion}
\label{sec:conclusion}

We present GBA-Facade, the first benchmark evaluating the GlobalBuildingAtlas for autonomous facade maintenance applications. By ingesting the complete 2.75B-building dataset into PostgreSQL and performing systematic quality auditing, we demonstrate that GBA enables planetary-scale market sizing for robotic facade cleaning, with a total addressable facade area of 4.18B~m\textsuperscript{2} across 408,812 high-rise buildings worldwide. Our 14-city benchmark reveals that the Pearl River Delta alone accounts for 58,742 buildings above 50m with 818M~m\textsuperscript{2} of facade surface. Our data quality audit reveals 65.9M negative-height artifacts (2.38\%) requiring correction, and validation against NYC PLUTO shows a height MAE of 3.12m ($R^2 = 0.84$). We release our complete pipeline to support further research at the intersection of earth observation data and robotics deployment.

% ─── Acknowledgments ──────────────────────────────────────────────────────
\section*{Acknowledgments}

We thank Zhu et al.\ at the Technical University of Munich for releasing the GlobalBuildingAtlas as open data. Compute infrastructure provided by AWS (eu-west-3). VitroBOT is supported by the French deeptech ecosystem.

% ─── References ───────────────────────────────────────────────────────────
\bibliographystyle{plainnat}

\begin{thebibliography}{10}

\bibitem[Zhu et~al.(2025)]{zhu2025gba}
X.~X. Zhu, S.~Chen, F.~Zhang, Y.~Shi, and Y.~Wang.
\newblock GlobalBuildingAtlas: an open global and complete dataset of building
  polygons, heights and LoD1 3D models.
\newblock \emph{Earth System Science Data}, 17(12):6647--6668, 2025.
\newblock \doi{10.5194/essd-17-6647-2025}.

\bibitem[Microsoft(2023)]{microsoft2023buildings}
Microsoft.
\newblock Global {ML} Building Footprints, 2023.
\newblock \url{https://github.com/microsoft/GlobalMLBuildingFootprints}.

\bibitem[Sirko et~al.(2023)]{sirko2023google}
W.~Sirko, E.~Marber, Y.~Ber, and J.~Quinn.
\newblock Google Open Buildings V3.
\newblock \emph{arXiv preprint arXiv:2310.11710}, 2023.

\bibitem[Ott et~al.(2020)]{ott2020facade}
C.~Ott, K.~Berns, and C.~L\"offler.
\newblock Autonomous facade cleaning robots: A survey of challenges and
  solutions.
\newblock \emph{Journal of Building Engineering}, 34:101935, 2020.

\bibitem[Wang et~al.(2021)]{wang2021bim}
J.~Wang, W.~Sun, J.~Li, and Z.~Zhu.
\newblock {BIM}-based facade cleaning robot path planning.
\newblock \emph{Automation in Construction}, 131:103882, 2021.

\bibitem[NYC~DCP(2023)]{nyc2023pluto}
New~York City Department of~City~Planning.
\newblock {PLUTO} and {MapPLUTO}, 2023.
\newblock \url{https://www.nyc.gov/site/planning/data-maps/open-data/dwn-pluto-mappluto.page}.

\bibitem[PostGIS(2024)]{postgis2024}
PostGIS Development Team.
\newblock PostGIS 3.4 documentation, 2024.
\newblock \url{https://postgis.net/documentation/}.

\end{thebibliography}

\end{document}
